% Fisica 1 - Fabrizio Cei - 09/02/2024
% Corpo Rigido - Pagina 21

\section{Moto di pura rotazione}

L'asse attorno a cui avviene è fisso, bloccato.

$d_i = d_i' = r_i \sin\theta_i = r_i' \sin\theta_i$ $\Rightarrow$ solo può traslare lungo l'asse di rotazione.

% Diagramma: corpo rigido con asse z fisso, masse m_i a distanza d_i dall'asse, forze F_i

\begin{itemize}
    \item $L_O$ = momento angolare assiale
    \item $I$ = momento di inerzia assiale
    \item $\tau_O$ = momento assiale della forza
\end{itemize}

È un ragionamento analogo a $L_O$, basta verificare che $r_i \sin\theta\, F_y = r_i' \sin\theta'\, F_y$.

\[
\begin{cases}
\vec{L} = I_z\,\vec{\omega} & \text{asse fisso principale} \\
L_z = I_z\,\omega & \text{asse fisso qualsiasi}
\end{cases}
\qquad
\begin{aligned}
& \text{accelerazione} \\
& \text{angolare}
\end{aligned}
\]

Studiamo la 2ª equazione cardinale per il corpo rigido che ruota attorno a $z$:

\paragraph{(1) Asse principale: $\vec{L}_{tot} = I_z\,\vec{\omega}$}

\begin{equation}
\vec{\tau}_{ext} = \frac{d\vec{L}_{tot}}{dt} \qquad \Rightarrow \qquad \boxed{\vec{\tau}_{ext}} = \frac{d(I_z\,\vec{\omega})}{dt} = I_z \cdot \frac{d(\vec{\omega})}{dt} = \boxed{I_z\,\vec{\alpha}}
\end{equation}

\paragraph{(2) Asse qualsiasi: $L_z = \omega I_z$}

\begin{equation}
\left[\tau_{ext}\right]_z = \frac{d(I_z\,\omega)}{dt} = \boxed{I_z\,\alpha}
\end{equation}

$\rightarrow$ $I_z$ è costante perché il corpo rigido è indeformabile e incomprimibile.

Ci sono però dei sistemi in cui $I_z$ può variare:

% Diagramma: sistema con masse che si spostano lungo un asse, cambiando la distribuzione di massa

\begin{equation}
\Rightarrow \tau_{ext,z} = \frac{d(L_z)}{dt} = I_z\,\alpha + \omega\,\frac{dI_z}{dt}
\end{equation}

$\hookrightarrow$ Questo caso è interessante se il sistema è \textbf{isolato} ($\tau_z = 0$):

\begin{align*}
I_z\,\alpha + \omega\,\frac{dI_z}{dt} = 0 \quad &\Rightarrow \quad \alpha = -\frac{d(I_z)}{dt}\cdot\frac{1}{I_z}\cdot\omega \\
&= \frac{d(\omega)}{dt}\cdot\frac{1}{\omega}
\end{align*}

da cui:
\begin{equation}
\frac{d(\omega)}{d(I_z)} = -\frac{\omega}{I_z}
\end{equation}

Cioè se $\omega$ aumenta ($d(\omega) > 0$) allora $I_z$ deve diminuire ($d(I_z) < 0$) e viceversa.

\begin{osservazione}{}
Le forze interne mutano il momento di inerzia.
\end{osservazione}
